% Copy-ready extension following the componentwise canonical atom spectrum.
% This fragment does not modify the authoritative manuscript automatically.

\subsection{Shared-point multiplicity spectra}
\label{shared-point-multiplicity-spectra}

Let $F$ be reduced and obligatory with at least one hyperedge, let
$A_1,\ldots,A_k$ be its canonical atoms, and let $c$ be its number of
connected components.  Write $P_{\rm sh}(F)$ for the points belonging to at
least two atoms and
\[
\mu(p)=|\{i:p\in V(A_i)\}|.
\]
The atom--shared-point incidence graph has the atoms and $P_{\rm sh}(F)$ as
its two classes.  By Theorem~\ref{theorem-canonical-atom-normal-form}, it is a
forest, and its components are precisely the components of $F$.

\begin{theorem}[Exact shared-point excess profile]
\label{theorem-exact-shared-point-excess-profile}
For every such $F$,
\[
\sum_{p\in P_{\rm sh}(F)}(\mu(p)-1)=k-c.
\]
Consequently the nonincreasing multiset
\[
\lambda(F)=\{\mu(p)-1:p\in P_{\rm sh}(F)\}
\]
is an integer partition of $k-c$.  Conversely, for any fixed list of $k$
canonical atom types and any $1\le c\le k$, every partition
$\lambda\vdash k-c$ occurs as the shared-point excess profile of a forest
assembly into $c$ components.  The assembly preserves all atom isomorphism
types.
\end{theorem}
\begin{proof}
If $t=|P_{\rm sh}(F)|$, the incidence forest has $k+t$ vertices, $c$
components, and therefore $k+t-c$ edges.  Summing degrees over the
shared-point class gives
\[
\sum_p\mu(p)=k+t-c,
\]
which proves the identity.

For the converse, reserve $c-1$ atoms as isolated components.  If
$\lambda=(\lambda_1,\ldots,\lambda_t)$, start with one of the remaining
atoms.  At step $j$, identify a fresh point of the current frontier atom with
fresh points of $\lambda_j$ unused atoms, and choose one of the new atoms as
the next frontier.  The resulting shared point has multiplicity
$\lambda_j+1$.  Since $\sum_j\lambda_j=k-c$, all remaining atoms are used.
Each atom is incident with at most two distinct shared points, so the
construction is valid for every canonical atom.  Its incidence graph is a
forest and the original atoms remain its canonical atoms.
\end{proof}

\begin{corollary}[Shared-point count and largest multiplicity]
\label{corollary-shared-point-count-largest-multiplicity}
Put $N=k-c$.  If $N=0$, there is no shared point.  If $N>0$ and $t$ shared
points are prescribed, then $1\le t\le N$, and the possible values of
$M=\max_p\mu(p)$ are exactly
\[
1+\left\lceil\frac Nt\right\rceil
\le M\le N-t+2.
\]
Every value in this interval occurs.  In particular, the number of
multiplicity profiles for fixed $k,c$ is the partition number $p(k-c)$.
\end{corollary}
\begin{proof}
A profile with $t$ shared points is a partition of $N$ into $t$ positive
parts.  Its largest part lies between $\lceil N/t\rceil$ and $N-t+1$.
Conversely, every integer in this interval is the largest part of such a
partition: after fixing the largest part $r$, distribute $N-r$ among the
remaining $t-1$ positive parts, each at most $r$.  Apply the theorem.
\end{proof}

\begin{proposition}[Capacity-aware incidence degrees and Prüfer count]
\label{proposition-capacity-aware-incidence-degrees}
Suppose $F$ is connected, $k\ge2$, and its $t$ shared points are labelled.
Let $d_i$ be the degree of atom $A_i$ in the incidence tree and let
$\mu_j$ be the degree of the $j$th shared point.  Then
\[
\sum_i(d_i-1)=t-1,
\qquad
\sum_j(\mu_j-1)=k-1.
\]
Conversely, let $1\le d_i\le |V(A_i)|$ and $2\le\mu_j$ satisfy
\[
\sum_i d_i=\sum_j\mu_j=k+t-1.
\]
There is a forest assembly with these incidence degrees.  For labelled atom
and point nodes, the number of incidence trees with the prescribed degree
sequences is
\[
\frac{(t-1)!(k-1)!}
 {\prod_i(d_i-1)!\prod_j(\mu_j-1)!}.
\]
\end{proposition}
\begin{proof}
The identities are the two degree sums in a bipartite tree.  Positive
bipartite degree sequences with common sum $k+t-1$ are realised by a tree;
this follows either by leaf deletion or from the bipartite Prüfer code.
Choose $d_i$ distinct vertices in atom $A_i$ and use them for its incident
shared-point nodes.  The tree condition prevents unintended identifications.
The Prüfer code consists of a word of length $t-1$ on the atom labels and a
word of length $k-1$ on the point labels.  Prescribing the occurrence counts
gives the displayed product of multinomial coefficients.
\end{proof}

Let $\mathcal K(s,\beta,c)$ be the exact atom-count set from
Proposition~\ref{proposition-componentwise-atom-count-spectrum}.  Combining
that proposition with the exact profile theorem gives
\[
\{\lambda(F):F\text{ has fixed }(s,\beta,c)\}
=
\bigcup_{k\in\mathcal K(s,\beta,c)}
\{\lambda:\lambda\vdash k-c\}.
\]
Thus the exact number of shared-point multiplicity profiles at fixed global
invariants is
\[
\sum_{k\in\mathcal K(s,\beta,c)}p(k-c).
\]
